\Subhead{2.1.3}{Упругие продольные и поперечные волны в твёрдом теле}

При малой деформации $e_{ij}=(\partial_i u_j+\partial_j u_i)/2$, $\sigma_{ij}=C_{ijkl}e_{kl}$ и $\rho\ddot u_i=\partial_j\sigma_{ij}$. Для плоской волны $\bm u=\bm a\e^{\ii(k\bm n\cdot\bm r-\omega t)}$ получаем уравнение Христоффеля
\begin{equation}
 \Gamma_{ik}(\bm n)a_k=\rho v^2a_i,\qquad
 \Gamma_{ik}=C_{ijkl}n_jn_l,\qquad v=\omega/k.
\end{equation}
Три собственных значения дают три скорости и поляризации. В изотропном теле
$\Gamma_{ik}=\mu\delta_{ik}+(\lambda+\mu)n_in_k$; следовательно,
\begin{equation}
 v_L=\sqrt{(\lambda+2\mu)/\rho},\qquad v_T=\sqrt{\mu/\rho}.
\end{equation}
Продольная поляризация параллельна $\bm n$, две поперечные — ортогональны ей и вырождены по скорости. В кристалле моды обычно квазипродольные и квазипоперечные; $\bm v_g=\nabla_{\bm k}\omega$ и фазовая скорость могут иметь разные направления.

Локальный баланс получается умножением уравнения движения на $\dot u_i$:
\begin{equation}
 w=\tfrac12\rho\dot{\bm u}^2+\tfrac12 C_{ijkl}e_{ij}e_{kl},\qquad
 S_j=-\sigma_{ij}\dot u_i,\qquad \partial_tw+\nabla\cdot\bm S=0.
\end{equation}
Для изотропной бегущей моды $\overline S=\rho v\omega^2|a|^2/2$. Это важное различие акустики жидкости и кристалла: твёрдое тело поддерживает сдвиговые напряжения и поперечные волны.
